\section{核心函数：\texttt{vmm\_map\_page}}

\codefile{src/kernel/mm/vmm.c}
\begin{lstlisting}[style=cstyle]
int vmm_map_page(uint32_t virt, uint32_t phys, uint32_t flags) {
    // 对齐检查
    if ((virt & 0xFFFu) || (phys & 0xFFFu)) {
        printk("[vmm] map: unaligned virt=0x%08x phys=0x%08x\n",
               virt, phys);
        return -1;
    }

    uint32_t pdi = pd_index(virt);
    uint32_t pti = pt_index(virt);

    page_table_t* pt = vmm_ensure_page_table(pdi);
    if (!pt) return -1;  // PMM OOM

    // 写入 PTE：高 20 位 = 物理帧地址, 低 12 位 = flags
    pt->entries[pti] = (phys & ~0xFFFu) | (flags & 0xFFFu);

    // 分页已开启则刷新 TLB
    if (g_vmm_ready) {
        vmm_invlpg(virt);
    }

    return 0;
}
\end{lstlisting}

\begin{figure}[H]
  \centering
  % figures/ch08/fig05-map-page-flow.tex — auto-extracted from chapter source
\begin{tikzpicture}[
  flowstep/.style={draw, rounded corners, minimum width=5cm, minimum height=0.7cm, font=\small},
  arr/.style={-{Stealth[length=3pt]}, thick},
]
  \node[flowstep, fill=gray!10]   (s1) at (0, 4)   {1. 对齐检查 virt \& phys};
  \node[flowstep, fill=red!10]    (s2) at (0, 3)   {2. pd\_index / pt\_index};
  \node[flowstep, fill=orange!10] (s3) at (0, 2)   {3. vmm\_ensure\_page\_table};
  \node[flowstep, fill=green!10]  (s4) at (0, 1)   {4. 写 PTE = phys | flags};
  \node[flowstep, fill=blue!10]   (s5) at (0, 0)   {5. invlpg（如果 paging 已开启）};

  \draw[arr] (s1) -- (s2);
  \draw[arr] (s2) -- (s3);
  \draw[arr] (s3) -- (s4);
  \draw[arr] (s4) -- (s5);
\end{tikzpicture}

  \caption{\texttt{vmm\_map\_page} 执行流程}
  \label{fig:map-page-flow}
\end{figure}

\begin{cpustate}
调用 \texttt{vmm\_map\_page(\hex{C0200000}, \hex{200000}, PTE\_PRESENT | PTE\_WRITABLE)} 后：

\begin{itemize}
  \item PDE[768] 的 Present 位被置 1（如果页表刚创建）
  \item PDE[768] 指向的 Page Table 的 PTE[512] = \hex{200003}（\hex{200000} | 0x3）
  \item 如果 TLB 中缓存了 \hex{C0200000} 的旧条目，\texttt{invlpg} 使其失效
  \item 此后访问虚拟地址 \hex{C0200000} 将读写物理地址 \hex{200000}
\end{itemize}
\end{cpustate}

% ============================================================================

